% EP04, EP08, EP10, EP11 + EP24. No core filtering is used on genuine V27.
% Complete editable argument relative to the terminal hidden Ramsey/reset
% and compatible tuple-readout interfaces in the internal V27/R18 source.
\section{The explicit unfiltered Ramsey expansion has expansion property}
\label{sec:ep-unfiltered}

Recall the \emph{expansion property} (EP) for a reasonable,
reduct-surjective precompact expansion
\(\mathcal K^*\to\mathcal K\): for every finite
\(A^*\in\mathcal K^*\) there is a finite
\(B\in\mathcal K\) such that \emph{every}
admissible expansion \(B^*\) of this fixed \(B\)
contains an embedding \(A^*\hookrightarrow B^*\).
This is a universal quantifier over all allowed
choices on \(B\), not merely existence of
one preferred homogeneous expansion.

We first record a finite-template criterion which is useful
outside CCO and, together with
\cref{thm:ep-automatic-core}, eliminates any need
for a separately selected minimal subflow or finite core retract
in our main construction.

\begin{theorem}[Simultaneous-singleton core criterion]
\label{thm:ep-finite-core-criterion}
Let \(\mathcal K\) be a hereditary finite complete
atomic binary age (allowing original named unary sorts).
Let \(\mathcal U_T\) be a reasonable, reduct-surjective,
pair-locally specified finite unary expansion with
a finite set \(T\) of realised complete unary labels:
\[
 (A,\lambda)\in\mathcal U_T
 \quad\Longleftrightarrow\quad
 A\in\mathcal K,\quad
 \col_A(x,y)=a\ \Longrightarrow\
         (\lambda(x),\lambda(y))\in R_a
 \quad(x\ne y).
 \tag{U1}\label{eq:ep-pair-local}
\]
Here every pair type is realised, the old atom
names and original sort predicates are retained,
and each \(R_a\subseteq T^2\) may overlap another.
Suppose \(\mathcal U_T\) has a reasonable hereditary
Ramsey expansion \(\mathcal R\) with JEP in which
there is exactly \emph{one expanded singleton type}
per label \(t\in T\). Then
\[
 \boxed{\quad
 \operatorname{EP}(\mathcal U_T/\mathcal K)
 \quad\Longleftrightarrow\quad
  \operatorname{End}(\mathbf T)
          =\operatorname{Aut}(\mathbf T).
 \quad}
 \tag{U2}\label{eq:ep-core-iff}
\]
No new unary label is individually named inside
the \emph{finite template} \(\mathbf T\); otherwise
the endomorphism test would be vacuous.
\end{theorem}
\begin{proof}
If a table endomorphism \(f:T\to T\) omits
some realised \(t\), any legal marking
\(\lambda\) of any old finite \(B\) can be
postcomposed with \(f\) by
\eqref{eq:ep-pair-local}. The result omits
\(t\), so it avoids the legal expanded singleton
of type \(t\). Thus EP fails.

For the converse fix \(A^*\in\mathcal U_T\).
For every \(\alpha\in\operatorname{Aut}(\mathbf T)\)
the inverse-relabelled object
\(\alpha^{-1}A^*\) is legal by
\eqref{eq:ep-pair-local}. Choose an
\(\mathcal R\)-expansion of each. Choose in
\(\mathcal R\) also a witness for \emph{every}
allowed old-colour/type pair
\((s,t)\in R_a\), including \(s=t\) on
\emph{distinct} real vertices, and one
singleton of every realised type. Finite
JEP embeds these finitely many objects into
one \(D\in\mathcal R\).

Enumerate the singleton shapes
\(s_1,\ldots,s_{|T|}\). Iterating the ordinary
Ramsey arrows for these shapes (with \(|T|\)
colours) produces \(C\in\mathcal R\)
such that every simultaneous choice of
\(|T|\)-colourings of the embeddings of the
\(s_i\) into \(C\) is separately constant on
some embedded copy of \(D\). For clarity,
one obtains this by choosing the arrows in
reverse order and peeling the resulting
embedded copies; it requires neither EP
nor any joint one-point language action.

Set \(B=C|_{\mathcal K}\). Given an arbitrary
legal competing \(T\)-labelling
\(\lambda':B\to T\), colour a reference
singleton \(s_i\) at vertex \(x\) by
\(\lambda'(x)\). On the selected \(D\)-copy
the competitor's label is a constant
\(f(t)\) on each \emph{reference} label \(t\).
All labels appear and the stored pair
witnesses force
\((f(s),f(t))\in R_a\) for every
\((s,t)\in R_a\). Thus \(f\) is an
endomorphism of \(\mathbf T\) preserving
the original sorts. It is an
automorphism \(\alpha\), by hypothesis.
The stored \(\alpha^{-1}A^*\)-copy
has exactly the prescribed competing
labels of \(A^*\). The single old \(B\)
was selected \emph{before} receiving
\(\lambda'\), as required.
\end{proof}

\begin{remark}[Finite normalisation as an independent fallback]
\label{rem:ep-core-filter}
For an \emph{arbitrary} pair-local expansion,
the image of a minimum-rank idempotent
endomorphism of \(\mathbf T\) is an induced
core. Restricting to its labels preserves
reasonableness, reduct-surjectivity and a
relational Ramsey provider, and yields EP
by \cref{thm:ep-finite-core-criterion}.
This is a general finite algorithm, not
part of our main V27 recipe. Its necessity
on the legitimate V27 family is ruled out
by \cref{thm:ep-automatic-core}. An
arbitrarily added free terminal unary
partition, or the nonmaximal safe
orientation in \cref{rem:ep-core-scope},
can have a noncore type table.
\end{remark}

\begin{lemma}[Symmetry of equivariant coordinates]
\label{lem:ep-unique-coordinate}
Let \(G\) act transitively on \(S\) and assume
all \(G\)-invariant binary relations on \(S\)
are symmetric. If \(\tau,\sigma:S\to B\)
are equivariant surjections onto the same
transitive \(G\)-set \(B\), then
\[
 \ker\tau=\ker\sigma,\qquad
 \sigma=f\circ\tau
\]
for an equivariant involution \(f:B\to B\).
If \(S\) has no proper finite-relative-index
invariant equivalence above \(\ker\tau\),
then \(f=\Id\) and \(\sigma=\tau\).
\end{lemma}
\begin{proof}
Define \(R(x,y)\) by
\(\tau(x)=\sigma(y)\). It is symmetric.
If \(\sigma(x)=\sigma(y)\), choose \(z\)
with \(\tau(z)=\sigma(x)\), use
\(R(z,x),R(z,y)\), and then symmetry
to get \(\tau(x)=\sigma(z)=\tau(y)\).
Interchanging \(\tau,\sigma\) gives equality
of kernels and an equivariant bijection
\(f\) with \(\sigma=f\tau\).
Symmetry of \(R\) also implies \(f^2=\Id\).
A nontrivial equivariant involution on a
transitive \(B\) is fixed-point-free, so
its two-element orbits pull back to an
invariant equivalence of relative index
two over \(\ker\tau\). Terminality excludes
this case.
\end{proof}

\begin{corollary}[A singleton-faithful terminal Ramsey provider]
\label{cor:ep-singleton-faithful}
In the functorially block-normalised,
finite-language V27 terminal hidden Ramsey
expansion, interpreted on the original real
roots by the full R18 unary/MI-tuple code,
there is exactly one expanded singleton
type for each realised full terminal unary
label.
\end{corollary}
\begin{proof}
Every unary hidden coordinate term on
one terminal transitive real sort gives
an equivariant surjection onto its named
target quotient sort. By
\cref{lem:ep-unique-coordinate}, two
terms targeting the same sort have
identical values. V27's
no-real-algebraic-transition theorem
rules out additional returning real
term values. Consequently one real
root has at most one unary term value
in each hidden target sort.

Order comparisons between values of
different target sorts are fixed by the
preassigned priority of hidden sort
symbols. Comparisons within one target
sort are equality for terms of the
same real root. All non-order unary
term/equality/cylinder data is invariant
under the transitive terminal
automorphism group and hence fixed by
its realised parent unary label. These
statements include any MR root interpreted
as a compatible MI tuple: each of its
coordinates is a unary term of a clone
representative and has the same
singleton normal form.
\end{proof}

The remaining ordering argument concerns the \emph{full fine
colours}, not merely the distributive lattice-ultrametric
equivalence reduct. It takes place inside the already
constructed terminal hidden Ramsey age and uses the
relative-reset semantic type of the V27 appendix.

\begin{lemma}[Finite holonomy vanishes in terminal quotients]
\label{lem:ep-terminal-holonomy}
An equivariant permutation of a transitive represented
terminal quotient sort is the identity. Hence orders
transported along definable quotient bijections have
no nontrivial closed-loop holonomy.
\end{lemma}
\begin{proof}
The graph of such a permutation is a binary invariant
orbital relation, and there are finitely many of them.
Its successive powers therefore eventually repeat,
so the permutation has finite order. A nonidentity
equivariant finite-order permutation is fixed-point-free
on each nontrivial orbit; these finite orbits give
a proper finite-relative-index invariant equivalence
above the kernel of the represented quotient map.
This contradicts terminality.
\end{proof}

\begin{lemma}[One complete generated type of an ordered reset edge]
\label{lem:ep-reset-complete-edge}
Let $J<J^+$ be one of the genuinely represented
terminal meet-irreducible cuts, let
$\mathfrak m\subseteq J^+\setminus J$
be ONE \emph{complete named} relative
$J$-semantic type, and suppose all
definable quotient-transition bijections
and the corresponding subquotient orders
are synchronized. Then every ordered pair
of $J$-classes $u<v$ with
$u\,\mathfrak m\,v$ has the same complete
finite unary-generated \emph{hidden expanded}
two-root type $H_{\mathfrak m}$, including
all unary-term equality codes, original
fine semantic cylinders and transported
order comparisons.

This statement is relative to V27's
\emph{full transition-enriched hidden
language}; it would not follow from the
bare equivalence reduct or merely from
the RAW equality-inclusive identity
\(\mathfrak m;\mathfrak m=J^+\) from
\cref{eq:typed-relative-reset-raw} without completeness
of the named semantic type.
\end{lemma}
\begin{proof}
The typed unary-coordinate normal form
\cref{lem:typed-coordinate-normal-form}
is stated for terms on a \emph{real} sort,
so we first justify its application here.
For a hidden unary term $\tau$ starting at
a proper quotient sort $B_J$, compose it
with the real projection $F_J:O\to B_J$.
The resulting real-source term
$\tau\circ F_J$ falls under that lemma.
Its equivalence kernel $L$ contains $J$,
since the term factors through $F_J$.
By the no-real-return theorem it is a
proper represented meet-irreducible,
not real equality. If $L=J$, then
the induced map from the $J$-quotient
is an equivariant bijection onto its
actual transitive target quotient.
Otherwise $L>J$ is a strict coarsening.
Because $J$ is meet-irreducible, it has the
unique cover $J^+$, and $J^+\le L$ in
the latter case. Thus for an
$\mathfrak m$-pair within one $J^+$-class
all strictly coarser term values at $u,v$
are equal at their corresponding transported
targets. Equality between two term
presentations, including those from
different roots, is determined by this
dichotomy and the complete $J$-semantic
pair type. On bijectively transported
$J$-coordinates the \emph{entire}
named $\mathfrak m$-cylinder fixes every
binary fine atomic relation between the
term values; the coarser values, being
unary images of one of the same two
root coordinates, introduce no new
independent pair choices.

Within each hidden target sort the
prescribed comparisons are either
equalities at a coarsening, or the
transported orientation of $u<v$ at
bijective $J$-level. Any two quotient
transition paths with the same target
agree, since terminality makes their
equivariant return holonomy trivial by
\cref{lem:ep-terminal-holonomy}. Between
different target sorts the hidden global
order comparison is the fixed priority
of the SORT SYMBOLS. Consequently every
unary-generated equality, fine-cylinder
and order code of the pair is forced by
the original named $\mathfrak m$ and
the orientation $u<v$.
\end{proof}

\begin{theorem}[Synchronized terminal subquotient-order EP]
\label{thm:ep-terminal-orders}
Let \(\mathcal H_{\mathrm{imag}}\) be the
functionally closed PURE-IMAGINARY hidden
age obtained from the terminal V27 hidden
envelope by retaining all proper represented
quotient sorts, their unary upward maps,
quotient transitions and COMPLETE named
fine semantic cylinders, but forgetting
the clone-real hidden sorts. On this age
consider the explicit order expansion:
for each represented meet-irreducible cut
\(J<J^+\), order \(J\)-classes inside
each \(J^+\)-class; transport each order
along its definable quotient-transition
orbit, choose one independent task per
orbit, and use a fixed hidden-sort priority
and functorial parent-coordinate
lexicographic reconstruction for all
remaining comparisons.

This is a reasonable, precompact Ramsey
order expansion with EP over the
\emph{full fine-coloured} pure-imaginary
hidden age \(\mathcal H_{\mathrm{imag}}\),
relative to the proved V27 typed hidden
Ramsey/SAP/reset source. It does NOT
claim relative EP for arbitrary orders
on deleted clone-real hidden vertices;
these orders are not part of the R18
MI/MR tuple readout.
\end{theorem}
\begin{proof}
The full terminal hidden class with arbitrary
global orders is hereditary Ramsey with JEP
and SAP by the internal V27 construction.
The class of proper imaginary coordinate
sorts is UNARY-CLOSED: by
\cref{lem:no-real-algebraic-transition}
no quotient-transition or upward unary map
from one proper imaginary sort has a real
target. Therefore
\cref{lem:foundational-ramsey-sort-restriction}
proves that the pure-imaginary ordered hidden
subage is itself Ramsey, JEP and SAP.
This argument uses actual induced closed
substructures, not an unsupported claim
that an arbitrary reduct of a Ramsey
class remains Ramsey.

Apply \cref{lem:foundational-ramsey-normalization} to the hierarchical
transition-synchronized normalizer of
\cref{lem:foundational-hierarchical-order-normalization}.
That normalizer is an idempotent endofunctor of the full
pure-imaginary ordered class, leaves vertex maps and unexpanded
hidden structures unchanged, and has exactly the desired normalized
objects as its fixed class. It already orders different sorts by the
fixed sort priority, so no separate block-normalization stage is needed.
The common lemma gives a hereditary Ramsey/JEP provider. This does not
assert that every partially ordered intermediate reduct is Ramsey.

The image expansion is also reasonable
and reduct-surjective: for any finite
unexpanded pure-imaginary \(B\), choose
an arbitrary full hidden order and
normalize it. Given an embedding of
an already normalized ordered \(A\)
into an unexpanded \(B\), first extend
the finite linear order of \(A\) to
an arbitrary total order of the full
underlying \(B\), then normalize.
Functoriality and idempotence give
\(N(B)|_A=N(A)=A\); hence the embedding
extends to a legal ordered expansion
of \(B\). Precompactness follows
from finiteness of the hidden sorts
and of the input structure.

Relative reversal of a SINGLE independent
transition-orbit task leaves all earlier
coarser task orders fixed; its effects
on finer parent-lex comparisons are
recomputed functorially by
\cref{lem:foundational-hierarchical-order-normalization}.
Consequently the intermediate reducts
used below genuinely differ by only
ONE task and its derived codes.

We verify the four actual hypotheses of the
abstract single-task criterion
\cref{thm:abstract-reset-order-EP}
for one orbit of a terminal represented
meet-irreducible cut $J<J^+$.
Earlier coarser task orders are held fixed as
part of the intermediate reduct; later finer
task orders are not fixed at this stage.
The \emph{full} transition-synchronized,
block-normalized PURE-IMAGINARY ordered
hidden class is the common Ramsey/JEP
provider throughout, so we never assume
an intermediate partially ordered class
is Ramsey.

First, the transition-synchronized orders
on $J$-classes inside each $J^+$-class
are free relative to the already fixed
coarser parent orders. Reversing this
one task \emph{simultaneously along all its
definable transition copies} therefore
produces another legal expansion over
the very same intermediate reduct.
Any dependent parent-tuple comparisons
are reconstructed by the fixed top-down
lexicographic rule, rather than being
independently prescribed. This is condition
(i) of the abstract criterion.

Next, the typed same-sort relative-reset theorem
\cref{thm:typed-relative-reset}, proved by
relativising the sparse-graph and private-theta
argument to the full SAP terminal typed source
(rather than merely carrying indices through
the historical one-sort proof), gives one
complete named symmetric $J$-semantic type
for the genuinely represented NONZERO cut
$J\in\mathcal I_a$:
\[
 \mathfrak m\subseteq J^+\setminus J,
 \quad \mathfrak m\starprod\mathfrak m=(J^+)^\circ,
 \quad \mathfrak m;\mathfrak m=J^+.
 \tag{O1}\label{eq:ep-reset}
\]
The first operation is the STRONG
typed diversity product and its
output is the ACTUAL diversity
palette \((J^+)^\circ\).
The final equality is the
equality-INCLUSIVE RAW relation
identity \cref{eq:typed-relative-reset-raw}.
We use the STRONG equality below
only on actual off-diagonal old
pair colours and never insert a
formal identity edge.
Every orientation of an $\mathfrak m$-edge
has the same \emph{entire} unary-generated
expanded two-root type
$H_{\mathfrak m}$ by
\cref{lem:ep-reset-complete-edge}.
This includes fine cylinders, equality,
upward transition terms and comparisons,
so condition (ii) holds.

For condition (iii), take an actual finite
functionally closed hidden object $A$
with $u<v$ two distinct $J$-classes inside
one $J^+$-class. We must realise the reset
\emph{over the actual complete hidden
two-root type of $(u,v)$}, not just over
some abstract pair of one $J^+$-class.

Embed the unexpanded $A$ in a full point-generated
coordinate expansion of the original terminal
fine-coloured homogeneous source, and choose
real representatives $x,y$ with
$F_J(x)=u$, $F_J(y)=v$. Their actual old
fine atomic pair colour $a$ belongs to
$J^+\setminus J$. The typed relative-reset theorem
\cref{thm:typed-relative-reset}
asserts the STRONG diversity equality
\(\mathfrak m\starprod\mathfrak m=(J^+)^\circ\),
in addition to its separately proved
RAW equality-inclusive counterpart.
We use the former, not a formal
identity edge, in this argument. Since
$a\in J^+\setminus J$ is a \emph{diversity}
atom, its inclusion in that strong product
is witnessed by a genuine old allowed
triangle on three distinct real vertices
with colours $(\mathfrak m,\mathfrak m,a)$.
By homogeneity on the old fine pair colour
$a$, transport that triangle's endpoints to
the \emph{specified} real pair $(x,y)$.
This supplies an actual real midpoint $z'$
whose $J$-class $z_0=F_J(z')$ satisfies
\[
          u\,\mathfrak m\,z_0\,\mathfrak m\,v.
\]
The ambient induced unary closure of $u,v$
already has EXACTLY the original hidden
fine-cylinder, equality and transition data
of $P=\langle u,v\rangle\le A$; we have
not replaced the prescribed pair by an
unrelated pair with merely the same
$J^+$-parent. This use of old pair
homogeneity is distinct from and prior to
the subsequent strong hidden amalgamation.
Because $\mathfrak m$ lies outside $J$,
$z_0$ differs from $u,v$ as a $J$-class.
Moreover $z_0$ cannot belong to the
unary-generated closure $\langle u,v\rangle$:
a unary term returning to the $J$-level
is a quotient-transition bijection, and
terminal no-holonomy makes every return
to the same named $J$ sort the identity.
All other terms from $u,v$ land in
strict coarsenings at least $J^+$,
whereas $z_0$ is a distinct $J$-point.
Thus $z_0$ is genuinely outside the
closed two-root base.

Take the finite unary closure
$T=\langle u,v,z_0\rangle$ with the same
induced complete two-root base
$P=\langle u,v\rangle$ as in $A$.
Choose the task-fibre order $u<z_0<v$;
it is legal because all three points
have the same strict $J^+$ parent and
the remaining comparisons are
functorially transported/extended.
Apply typed hidden SAP
\cref{thm:typed-hidden-sap} to $T$ and
$A$ over the \emph{complete closed} $P$.
The two given orders have identical
restrictions to $P$. On each represented
hidden sort and each allowed task fibre,
amalgamate these finite linear orders
over their common suborder by shuffling
the off-base points in each interval
between consecutive base points.
For the current $J$ task this keeps
$u<z_0<v$; other existing comparisons
in $A$ are unchanged. Recompute all
functorial parent-lex and transition
comparison codes from those joint
orders. This is the joint-order construction just described, using
\cref{thm:typed-hidden-sap} and the earlier sort-block normalization.
The resulting copy $z$ is fresh over
$A$ by strongness. This is exactly
\cref{lem:abstract-fresh-reset-midpoint};
repeat the construction for every
comparison in a finite task-fibre
configuration. No extension of a
prescribed old fine-colour matrix by
guessing a new colour is used:
the genuine hidden SAP theorem
constructs the common object.

The complete ordered PURE-IMAGINARY
hidden class is Ramsey and has JEP by
the V27 hidden Ramsey source, the
unary-closed sort restriction and
functorial hierarchical normalization
already checked at the start of this proof.
Consequently the abstract midpoint
theorem
\cref{thm:abstract-reset-order-EP}
gives EP for this \emph{one} current
task over all earlier tasks, with an old
host chosen before any competing order.
Its proof stores the task target and
its reversal, inserts the two reset
edges at every comparison, applies one
\(
 C\longrightarrow(D)^{H_{\mathfrak m}}_2
\)
arrow in the SAME full hidden Ramsey
class, and uses the constant comparison
sign to select the correct target.
The source theorem therefore does not
need a Ramsey statement for partially
ordered intermediate structures.

Finally process the finitely many
independent transition-orbits of cuts
from coarser to finer. Their relative
EP properties compose by
\cref{lem:foundational-relative-EP-composition},
whose proof transports \emph{arbitrary}
expanded hosts along actual embeddings
and preserves the quantifier order.
All the other parent-lex comparisons
are functorially determined by the
task orders and fixed sort priority,
so the final expanded order class
has EP over the full fine-coloured
PURE-IMAGINARY terminal hidden reduct.
No order on a deleted clone-real vertex
was used in the relative EP quantifiers.
\end{proof}

\begin{lemma}[Cofinal finite MI-tuple EP transfer]
\label{lem:ep-tuple-transfer}
Let \(\mathcal H^o\to\mathcal H\) be a reasonable
hidden finite-relational/unary-functional
precompact expansion with EP.
Suppose an interpretation \(N(D)\)
reads visible roots as finite compatible
tuples of selected unary hidden
quotient coordinates, and \(h(A)\)
is their finite unary-generated support.
Assume:
\begin{enumerate}[label=(\textup{\roman*}),leftmargin=*]
\item the unexpanded \(h(A)\) depends only
on the old fine-coloured visible \(A\);
\item every finite hidden \(D\) extends
to a finite \(E\) with \(h(N(E))=E\);
\item \emph{every} allowed visible
expansion of \(N(E)\) lifts through
its full comparison/equality code
to a hidden \(\mathcal H^o\)
expansion of the \emph{same} \(E\);
\item visible embeddings extend uniquely
to generated supports, and \(N\)
recovers every original fine colour.
\end{enumerate}
Then the visible expansion has EP.
The full R18 compatible MI-tuple
readout of Part~III satisfies these
four contracts, including the MR case.
\end{lemma}
\begin{proof}
The abstract transfer under (i)--(iv) is already
the general Part~I theorem
\cref{thm:foundational-cofinal-tuple-EP}.
Here it remains to audit the four
contracts for the actual compatible
R18 MI/MR tuple interpretation.
We verify the hypotheses of the two general
Foundations lemmas for the \emph{actual} R18
terminal hidden source. First delete the
clone-real hidden sorts. No unary transition
has its target in those real sorts, by
\cref{lem:no-real-algebraic-transition}, so
the pure-imaginary hidden reduct is
functionally closed and its finite age
is represented in the pure-imaginary part
of the actual terminal Fra\"iss\'e source.

Every proper imaginary point $b$ in this
actual ambient hidden structure is an
original represented coordinate or a
quotient-transition value of such a
coordinate. Choose a real representative
$x_b$ of its corresponding quotient class.
By the typed coordinate-term normal form
\cref{lem:typed-coordinate-normal-form},
the kernel of this coordinate term is a
proper represented meet-irreducible, and
the value is obtained from that coordinate
by unary quotient maps/transitions.
If the target coordinate $J$ is one of
the selected $J_{s,i}$, it is a coordinate
of the root tuple $\eta_s(x_b)$.
For a represented proper meet-irreducible
coordinate $J$ with $K_s=\bigcap_iJ_{s,i}\le J$,
the explicit finite-lattice cofinality
\cref{lem:foundational-MI-coordinate-cofinal}
gives a selected $J_{s,i}\le J$, so its
value is an upward unary image of
$\eta_s(x_b)$. More generally the
transition normal form expresses any
represented proper target by this finite
unary closure. The compatibility
predicate for $\eta_s(x_b)$ is the
quantifier-free R18 predicate
\cref{lem:ramsey-mi-domain}.

Consequently \emph{every} point of the
ambient pure-imaginary hidden source
is a unary-term value of a valid visible
MI tuple. The cofinal root-coverage lemma
\cref{lem:foundational-root-cover-cofinal}
now proves (ii): for each finite
pure-imaginary $D$ it constructs a finite
$E\supseteq D$ from the unary closure
of finitely many chosen real-representative
root tuples, with $h(N(E))=E$. This also
handles orphan coordinates and their
\emph{missing packet partners}, since
the entire unary closure of each root
tuple is included in $E$.

For (i), the unexpanded visible structure
at this relative EP stage contains the
\emph{terminal unary labels} and all
original fine binary atoms. The coordinate
terms and proper quotient transports are
canonical and $0$-definable from that
source: every unexpanded term equality,
fine-cylinder atom and unary incidence is
a union of its complete atomic old types,
as proved in \cref{cor:ramsey-small-language}.
Thus the unexpanded generated support
$h(B)$ is canonical over the visible
roots, not an arbitrary additional
imaginary coordinate assignment.
For the full-code EP step we keep ALL
unary terms, across-root equalities,
fine-cylinder names and comparisons
within each hidden target sort.

Now take the fixed $E$ supplied by
(ii), so $h(N(E))=E$, and an
\emph{arbitrary} competing allowed
full-code visible expansion $B^o$ of
$B=N(E)|_{\mathcal V}$. Allowedness means
that $B^o$ comes from an induced full
code on some hidden expanded object,
not necessarily on $E$.
The common same-host lift lemma
\cref{lem:foundational-same-host-code}
transports its complete generated
hidden code to the \emph{existing}
unexpanded $E$ using (i). This
produces one allowed hidden expansion
$E^o$ on exactly the same underlying
$E$ and realizes $B^o$, establishing
(iii) without selecting a convenient
alternative hidden host.

The all-embeddings identity of
\cref{lem:foundational-full-code-embeddings}
gives (iv). This argument does not
infer EP by quotienting an arbitrary
Ramsey expansion or by silently
discarding missing-partner
comparison information.
\end{proof}

\begin{theorem}[Explicit Ramsey and expansion-property theorem]
\label{thm:ep-unfiltered-main}
For every one-original-sort countable transitive
homogeneous complete finite symmetric binary
AP+homomorphic-CCO source, the following
\emph{unfiltered} finite unary/binary
expansion is reasonable, reduct-surjective,
precompact and Ramsey, and has the
\emph{expansion property} over the
original age:
\begin{enumerate}[label=(\textup{\arabic*}),leftmargin=*]
\item pass to the pure-clone SAP cover and
perform any \emph{genuine} finite
maximal-exact-character safe terminalisation;
retain \emph{every} realised joint
terminal unary type, removing \emph{none};
\item on terminal hidden MI quotients
use transition-synchronised
subquotient orders, fixed sort-block
priorities and parent-lex reconstruction;
\item read them back on the
original vertices as compatible
MI-coordinate tuples with the
finite \emph{full} unary/binary
term-equality, cylinder and
comparison code.
\end{enumerate}
No canonical choice of tower or
minimality of the signature is claimed.
\end{theorem}
\begin{proof}
R17's clone stability shows the
genuine safe terminal labels are
constant on full clone fibres,
so they label the original real
points. Its pair-locality is the
realised full groupoid-character
pair rule \eqref{eq:ep-pair-local};
the clone atom has type support
equal to the diagonal and adds
no endomorphism restriction.

The unfiltered table \(\mathbf T\)
is core by \cref{thm:ep-automatic-core}.
The block-normalised hidden Ramsey
provider has one expanded singleton
type per full terminal label by
\cref{cor:ep-singleton-faithful}.
Apply \cref{thm:ep-finite-core-criterion}:
the \emph{unfiltered} unary phase age
has EP over the original age.
The order expansion has relative EP by
\cref{thm:ep-terminal-orders};
the interpreted visible order age
has relative EP over the unary phase
age by \cref{lem:ep-tuple-transfer}.
These two relative EP statements
compose with unary EP by the
precise finite-chain quantifier lemma
\cref{lem:foundational-relative-EP-composition}.
First choose the necessary host at the
most expanded intermediate stage, then
successively choose hosts for its reducts.
For EVERY competing expansion of the
ONE final unexpanded host, restricting
its intermediate reduct along the
guaranteed embeddings supplies the
required preceding expanded witnesses.
The composition lemma uses heredity
and reduct-surjectivity; JEP and
precompactness are NOT needed for
this quantifier step. No arbitrary
order choice or phase marking is
selected after the final original
host is fixed.

Ramsey, reasonableness and
precompactness of the compatible
tuple readout were already
established in
\cref{thm:ramsey-general-ap,thm:ramsey-tuple-readout}; the
functorial order restriction
preserves their Ramsey arrows.
All displayed comparison codes
are unary or binary and there are
only finitely many uniform hidden
unary term normal forms.
Thus the entire visible expansion
is finite-language and at most
binary. EP of every layer has
now been verified rather than
inferred from its Ramsey property.
\end{proof}

\begin{remark}[Precise status]
\label{rem:ep-final-audit}
The core/label theorem is a direct
diversity-rank deduction of
\cref{sec:ep-exact-traces,sec:ep-automatic-core}
from the represented exact, quotient
and safe-character contracts.
The order and tuple clauses above
are direct finite categorical
arguments using substantial
terminal V27 hidden SAP/relative-reset
and R18 full-code inputs.
Those in-house source proofs are
reproduced or connected in the
technical appendices, but their
end-to-end independent line audit
and Lean formalisation are
\emph{not} claimed. The distinction
between original homomorphic CCO
and any stronger injective completion
statement is unchanged.
\end{remark}
